\scanpage{13}
Cas des groupes BT tronqués annulés par $p$.

Signifie que $F_{\mathcal M},V_{\mathcal M}$ forment suite exacte
\[
\Longleftrightarrow F_M,V_M\quad\text{\unsure{forme|forment} suite exacte}\;\ldots\ldots
\]
Alors la donnée b) implique, passant au $\underline H^1$
\margin{$\left[\begin{aligned}
\omega&=\underline H^1(\ell_\bullet[-1])\\
&=\underline H^0(\ell_\bullet)\\
\nu'&=\underline H^1(\check\ell'_\bullet)
\end{aligned}\right]$\par
$(\omega,\nu'\text{ loc. libres})$}
$M\inserted{{}=\varepsilon^*(\mathcal M)}$ filtré par
\[
0\mathbin{-}\omega\mathbin{-}M\mathbin{-}\nu'\longrightarrow0
\]
de façon que dans
\diagram{\begin{tikzcd}[column sep=small]
0\arrow[r,dash]&\varphi^*(\omega)\arrow[r]&\varphi^*(M)\arrow[r,dash]\arrow[d,equal]&\varphi^*(\nu')\arrow[r,dash]&0,\\
&&\mathcal M^{(p)}&&
\end{tikzcd}}
$\varphi^*(\omega)$ ne soit autre que $\operatorname{Im}(\mathcal M\xrightarrow{V_{\mathcal M}}\mathcal M^{(p)})$.
\margin{\uline{Cela détermine}\par\uline{$\omega$ de façon}\par\uuline{unique !}}

De $=$, passant au $\underline H^0$, on trouve
\margin{$\left[\begin{aligned}
\unsure{n|\nu}&=\underline H^0(\ell_\bullet[-1])\\
&=\underline H^{-1}(\ell_\bullet)\\
t'&=\underline H^0(\check\ell'_\bullet)
\end{aligned}\right]$}
\[
0\mathbin{-}\unsure{n|\nu}\mathbin{-}M\mathbin{-}t'\mathbin{-}0
\]
de façon que dans
\[
0\mathbin{-}\varphi^*(\unsure{n|\nu})\mathbin{-}\varphi^*(M)\mathbin{-}\varphi^*(t')\mathbin{-}0
\]
on ait
\[
\varphi^*(\unsure{n|\nu})=\operatorname{Ker}(\mathcal M^{(p)}\overset{F_{\mathcal M}}{\mathbin{-}}\mathcal M)
\]
Comme $\operatorname{Ker}F_{\mathcal M}=\operatorname{Im}V_{\mathcal M}$, on en conclut
\[
\varphi^*(\unsure{n|\nu})=\varphi^*(\omega)\qquad\text{d'où }\boxed{\unsure{n|\nu}=\omega},\quad\boxed{t'=\nu'}
\]
La seule donnée des \deleted{\unsure{suffit}}\inserted{$\mathcal M,F_{\mathcal M},V_{\mathcal M}\ (\Rightarrow\omega\unsure{\subset}M=\varphi^*(\mathcal M))$} permet elle de récupérer toute la donnée b) ? On voit qu'on pourra construire
\[
\omega\xrightarrow{V_\omega}\omega^{(p)},\qquad t'^{(p)}\xrightarrow{F_{t'}}t'
\]
\margin{$(\omega'=\check t',\quad\omega'\overset{V_{\omega'}}{\underset{\Vert\,{}^tF_{t'}}{\mathbin{-}}}\omega'^{(p)})$}
Car:
\[
\omega^{(p)}=\operatorname{Im}(M\deleted{^{\unsure{p}}}\overset{V_M}{\mathbin{-}}M^{(p)}),
\]
donc $V_M|\omega$ applique $\omega$ dans $\omega^{(p)}$
\[
t'^{(p)}=\operatorname{Coker}(M\xrightarrow{V_M}M^{(p)}),
\]
donc $F_M:M^{(p)}\to M$ se factorise par $t'^{(p)}\unsure{\mathbin{-}|\hookrightarrow}M$, d'où $F_{t'}:t'^{(p)}\to M\to t'$.
\margin{\deleted{\diagram{\begin{tikzcd}[column sep=small,row sep=small]
M\arrow[r]\arrow[d,"V_M"']\arrow[dr]&\omega\\
M^{(p)}\arrow[r,dash]&\omega^{(p)}\arrow[u,"F_\omega"']
\end{tikzcd}}}
\deleted{\diagram{\begin{tikzcd}[column sep=small,row sep=small]
M\arrow[r]\arrow[d]&t\\
M^{(p)}\arrow[r]&\unsure{t^{(p)}}
\end{tikzcd}}}}
\mspage{12}
